\clearpage
\fancyhead[L]{\small TAVOLE DELLE PARALLASSI.}
\hypertarget{native-604-lower}{}
\begin{center}\small Göttingische gelehrte Anzeigen. Stück 108. Seite 1073..1077. 1817 Juli 7.\\[3pt]\rule{.72\linewidth}{.3pt}\end{center}

\emph{Tavole delle parallassi di altezza, di longitudine e di latitudine calcolate dagli astronomi dell' osservatorio dell' università Gregoriana nel collegio romano.} Rom 1816. Nella stamperia de Romanis. XXXIV und 122 Seiten in Folio.

Die Bedeckungen der Fixsterne vom Monde und die Sonnenfinsternisse sind für die Astronomie und Geographie von so grosser Wichtigkeit, dass jeder Beitrag zur Erleichterung der darauf Bezug habenden Rechnungen mit Dank aufgenommen zu werden verdient. Die sogenannte Methode des Neunzigsten ist noch immer bei den Astronomen zur Berechnung jener Phänomene am meisten im Gebrauch, und die Berechnung der Längen- und Breitenparallaxe des Mondes macht einen Haupttheil derselben aus. Vorliegende Tafeln sind dazu bestimmt, diese Berechnung der Parallaxen abzukürzen. Der Plan dazu rührt von \textsc{Andreas Conti} her, einem der Astronomen der Römischen Sternwarte; bei der Ausführung dieser weitläufigen Arbeit wurde er von seinen Collegen \textsc{Joseph Calandrelli} und \textsc{Jacob Reichenbach} unterstützt.

Um unsere Leser in den Stand zu setzen, sich von der Einrichtung und dem Gebrauche dieser Tafeln, und von dem dadurch zu erreichenden Gewinn eine Vorstellung zu machen, müssen wir sie einzeln näher zergliedern. Der Tafeln sind eigentlich \emph{drei}; die erste enthält 93, die zweite und dritte jede 12 Seiten. Dazu kommen noch auf einem besondern Blatt einige kleine Hülfstafeln. Die erste Tafel enthält nichts anderes, als die Producte der halben Horizontalparallaxe in die Cosinus aller Winkel. Sie hat also doppelte Eingänge. Das eine Argument ist die Horizontalparallaxe, deren 50 Werthe von $53'\,20''$ bis $61'\,30''$ von 10 zu 10 Secunden fortlaufen. Die sämmtlichen Winkel des Quadranten von 10 zu 10 Minuten sind das zweite Argument. Jede Seite ist gespalten; die grössere Hälfte enthält die eben genannten Producte, die andere kleinere Hälfte gibt unter dem Titel, Supplement der ersten Tafel, die Producte derselben Cosinus in die halben Horizontalparallaxen der Planeten (von $1''$ bis $40''$); und ausserdem nochmals die Producte dieser Cosinus in alle einzelnen Minuten von $1'$ bis $9'$, und in die Zehner von Minuten von $10'$ bis $90'$. Der Raum ist also, wie man sieht, nicht gespart.

\SourceLeaf{605}Die zweite Tafel ist gleichfalls mit doppeltem Eingange; sie gibt uns den Werth des Products
$\alpha(\sec L - 1)$ für die Argumente $L$ und $\alpha$, jenes von 0 bis $6^{\circ}$ durch alle Zehner von Minuten, dieses von
0 bis $60'\,30''$ von halber zu halber Minute, und dann noch von $1''$ bis $35''$ durch alle ungeraden Secunden
genommen.

Die dritte Tafel enthält 12 Fixsterne, die vom Monde häufig bedeckt werden, den Werth eines
Hülfswinkel $\varphi$, welcher theils von der Breite des Fixsterns, theils von einem unten zu erklärenden
Winkel $D - \tfrac{1}{2}\Pi$ abhängig ist: die Werthe dieses Winkels laufen durch den ganzen Quadranten von
20 zu 2c Minuten.

Die drei kleinen Hülfstafeln enthalten noch: den Unterschied zwischen der Tangente und dem
Bogen, Verwandlung der Secunden in Decimaltheile der Minute und den Werth des Products,
$\operatorname{tang} \tfrac{1}{2}(D+\tfrac{1}{2}\Pi)\cdot\sin L'$, wo das erste Argument $D+\tfrac{1}{2}\Pi$ alle Winkel des Quadranten, das zweite $L'$
die Winkel von 0 bis $6^{\circ}$ von Grad zu Grad umfasst.

Dies ist der Inhalt der Tafeln; es bleibt uns noch übrig, den Gebrauch, welcher davon gemacht
werden soll, zu erklären.

Die Höhenparallaxe erhält man sofort durch Verdoppelung dessen, was die erste Tafel gibt, wenn
man in dieselbe mit der Horizontalparallaxe und der scheinbaren Höhe eingeht. Ist nicht diese, sondern die wahre Höhe gegeben, so wird man das Gesuchte durch wiederholte Annäherung erhalten; dieselbe Bemerkung gilt auch für das Folgende, in so fern die anzuwendenden Argumente nicht unmittelbar gegeben sind.

Die Parallaxe der Länge $\Pi$ wird, immer hinreichend genau, durch die Formel
\[
\Pi = \dfrac{P \sin h \sin(D+\Pi)}{\cos L}
\]
bestimmt, wo $P$ die Horizontalparallaxe, $h$ die Höhe des Neunzigsten, $D$ wahre Länge weniger Länge
des Neunzigsten, $L$ die wahre Breite bedeutet. Ist $\alpha$ die Differenz der Resultate der ersten Tafel, wenn
man in dieselbe einmal mit $P$ und $D-h+\Pi$, und dann mit $P$ und $D+h+\Pi$ eingeht, so wird $\Pi$
gleich sein dem Aggregat von $\alpha$ und dem Resultate der zweiten Tafel, wenn man in diese mit den Argumenten $\alpha$ und $L$ eingeht.

Die Breitenparallaxe $\pi$ ergibt sich hinreichend genau aus der Formel:
\[
\pi = -P[\cos h \cos(L+\pi) - \sin h \sin(L+\pi)\cos(D+\tfrac{1}{2}\Pi)]
\]
oder aus dieser
\[
\pi = -\dfrac{P\cos(h+\varphi)\cos(L+\pi)}{\cos\varphi}
\]
wo $\varphi$ einen durch die Formel
\[
\operatorname{tang} \varphi = \operatorname{tang}(L+\pi)\cdot\cos(D+\tfrac{1}{2}\Pi)
\]
bestimmten Hülfswinkel bedeutet. Da die letzte Formel für $\pi$ dieselbe Gestalt hat, wie die Formel für
$\Pi$, so sieht man leicht, wie zu ihrer Berechnung die erste und zweite Tafel angewandt werden können.

Die Hülfsmittel $\varphi$ gibt die dritte Tafel, wenn die Bedeckung eines der zwölf Sterne derselben
nach \textsc{Carlini}'s Methode berechnet werden soll; allgemein aber findet man ihn durch das Supplement der
ersten Tafel verbunden mit der ersten Hülfstafel. Da der Hülfswinkel $\varphi$ von der scheinbaren Breite
abhängt, so wird, wenn nicht diese, sondern die wahre Breite gegeben ist, folgendes Verfahren vorgeschrieben. Setzt man wie oben:
\[
P \sin h \sin(D+\Pi) = \alpha
\]
und zugleich
\SourceLeaf{606}\[
\sin(L+\pi)\,\operatorname{tang}\tfrac{1}{2}(D+\tfrac{1}{2}\Pi) = M
\]
so ist, hinreichend genau,
\[
\pi = -P\cos(h+\lambda+\pi) - \alpha M
\]
wo $M$ durch die letzte Hülfstafel, und $\alpha$ und $P\cos(h+\lambda+\pi)$ mit Hülfe der ersten Tafel gefunden werden. Wiederholte Annäherung ist auch hier unvermeidlich.

Sollen wir nun unser Urtheil über diese Tafeln offen erklären, so müssen wir gestehen, dass wir eine Erleichterung der Parallaxenrechnungen durch dieselben nicht finden können, sondern die Rechnung mit den gewöhnlichen trigonometrischen Tafeln zum wenigsten für eben so bequem halten, als die Anwendung dieser Specialtafeln. Druck und Papier sind übrigens schön, und so, wie man es allen Zahlenwerken wünschen möchte.

Merkwürdig ist noch die Zueignung dieses Werks an den Papst. Man freut sich über die Unterstützung, welche dieser der Römischen Sternwarte angedeihen lässt. Aber was soll man denken von dem leidenschaftlichen Eifern der Römischen Astronomen gegen `die Gottlosen, welche von der Astronomie den schändlichsten und schrecklichsten Missbrauch machen, das Herz verderben, die wahren Gläubigen verhöhnen und als unsinnige Anbeter der Sonne verspotten.' Wer versteht diese seltsamen Beschuldigungen, und wer würde, ohne eine unzweideutige Hindeutung in dieser Zueignung, errathen haben, dass sie einem grossen Geometer gelten sollen, welcher gewagt hat, eine Hypothese zur physischen Erklärung der Bildung des Sonnensystems aufzustellen?
